\pdfoutput=1

\documentclass[11pt]{article}

\usepackage[margin=1in]{geometry}
\usepackage{amsmath,amssymb,amsthm}
\usepackage{booktabs}
\usepackage{graphicx}
\usepackage{hyperref}
\usepackage{xcolor}
\usepackage{array}
\usepackage{enumitem}
\usepackage{multirow}
\usepackage{float}
\usepackage{lineno}
\usepackage[numbers,sort&compress]{natbib}
\graphicspath{{figures/}}

\newcommand{\IIA}{\mathrm{IIA}}
\newcommand{\R}{\mathbb{R}}
\newcommand{\KL}{D_\mathrm{KL}}
\newcommand{\zc}{z_{\text{causal}}}
\newcommand{\zn}{z_{\text{nuisance}}}
\newcommand{\TODO}[1]{\textcolor{red}{[#1]}}

\renewcommand{\arraystretch}{1.2}

\title{Interchange Accuracy Does Not Separate Constrained From Unconstrained
Alignment Maps; Intervention Diversity and Distribution Shift Do}

\author{Anonymous}
\date{}

\begin{document}
\linenumbers
\maketitle

\begin{abstract}
\citet{sutter2025nonlinear} showed that unconstrained nonlinear alignment maps
make causal abstraction uninformative, reaching 100\% interchange intervention
accuracy (IIA) on randomly initialised language models. We compare two nonlinear
alignment maps trained on the same interchange objective, differing in whether
that objective is optimised subject to a variational autoencoder's reconstruction
and regularisation terms. On indirect object identification in GPT-2 at $k = 1$
both reach $\IIA = 1.000$, so accuracy does not distinguish them. Two other
measurements do. An intervention diversity ratio, comparing the spread of
intervened activations against natural ones, is $0.05$ for the unconstrained map
and $0.83$ for the constrained map. Under distribution shift the two diverge
further: on held-out entities the unconstrained map reports $\IIA = 1.000$ and
the constrained map $0.001$, while on held-out templates the constrained map
reports $0.993$ against $0.537$. On grokked modular addition across subspace
dimensions 1 to 32, the unconstrained map does not exceed $\IIA = 0.861$, linear
DAS reaches $1.000$ at $k = 8$, and the constrained map reaches $1.000$ at
$k = 1$. We do not establish that either map recovers the model's causal
variable, and the reading of the held-out-entity result rests on an assumption
about the task that we state rather than test. The supported claim is narrower:
accuracy alone did not separate these two maps, and two inexpensive additional
measurements did.
\end{abstract}

\section{Introduction}
\label{sec:intro}

Causal abstraction formalises what it means for a network to implement an
algorithm: a model variable is causal if interventions on its internal
representation behave like interventions on a high-level algorithmic variable
\citep{geiger2021causal}. Distributed Alignment Search (DAS) operationalises this
by learning a rotation that identifies a $k$-dimensional subspace of the residual
stream, swapping the in-subspace component between two inputs, and measuring
whether the model produces the counterfactual output \citep{geiger2023finding}.
The fraction of successful swaps is the interchange intervention accuracy.

DAS constrains the alignment map to be linear. \citet{sutter2025nonlinear} showed
that removing this constraint makes the framework uninformative: sufficiently
expressive maps can align any network to any algorithm, and in their experiments
randomly initialised language models reached 100\% IIA on indirect object
identification. High accuracy obtained this way says nothing about the network.

One mechanism for this is visible in the training setup. A map trained on
interchange loss alone has no term requiring its encoder and decoder to be
approximate inverses. An optimiser could encode the label in one coordinate and
have the decoder emit a fixed activation per label, discarding the base input,
satisfying the interchange criterion without using the model's computation. We
refer to this as a lookup-table solution. We do not prove that unconstrained maps
converge to it; the measurements below are consistent with it rather than
demonstrations of it.

This paper compares two nonlinear alignment maps that optimise the same
interchange objective and differ in what constrains it. The unconstrained map is
a multilayer perceptron featuriser trained on interchange loss alone. The
constrained map adds the interchange term to a variational autoencoder's
reconstruction, regularisation, and supervision terms, so both maximise
interchange accuracy on identical data. Holding the objective fixed isolates the
constraints as a group, though not any single constraint; \S\ref{sec:ablation}
addresses that separately.

Because accuracy cannot by itself distinguish a lookup table from other
solutions, we report two further measurements alongside it: a diversity ratio
testing whether the decoder's output varies with the base input, and behaviour on
two distribution shifts, one of which we expect to preserve the causal variable
and one of which we do not.

\section{Results}
\label{sec:results}

\subsection{Accuracy and diversity}
\label{sec:vacuity}

If a decoder returns the same activation regardless of the base input, then
intervening on different bases with a shared source produces identical outputs.
The intervention diversity ratio $\rho$ (Methods) compares the spread of
intervened activations against the spread of natural ones. Values near zero
indicate that intervened activations for a given source label have collapsed
together; values near one indicate they retain the variation seen in natural
activations. We use $\rho$ as a proxy for whether the base input is being used,
and note that it is a summary statistic we introduce rather than a validated
measure of faithfulness.

\begin{table}[t]
\centering
\small
\setlength{\tabcolsep}{4pt}
\begin{tabular}{lc rrrr rr}
\toprule
Task & $k$ & DAS & NL-DAS & NL-DAS+r & Structured VAE & $\rho_\text{NL}$ & $\rho_\text{VAE}$ \\
\midrule
IOI & 1 & 0.183 & 1.000 & 0.867 & 1.000 & 0.05 & 0.83 \\
IOI & 2 & 0.322 & 1.000 & 0.700 & 0.989 & 0.05 & 0.89 \\
IOI & 4 & 0.572 & 1.000 & 0.811 & 0.989 & 0.05 & 0.92 \\
\midrule
Addition & 1 & 0.006 & 0.083 & 0.028 & 1.000 & 0.12 & 1.00 \\
Addition & 2 & 0.022 & 0.422 & 0.039 & 1.000 & 0.29 & 1.00 \\
Addition & 4 & 0.283 & 0.689 & 0.033 & 1.000 & 0.43 & 1.00 \\
\midrule
Mult.\ & 1 & 0.028 & 0.239 & 0.028 & 1.000 & 0.65 & 1.00 \\
Mult.\ & 2 & 0.072 & 0.544 & 0.111 & 1.000 & 1.00 & 1.00 \\
Mult.\ & 4 & 0.283 & 0.933 & 0.372 & 1.000 & 1.12 & 0.99 \\
\bottomrule
\end{tabular}
\caption{Interchange accuracy and diversity ratio. On indirect object
identification both nonlinear maps reach $\IIA = 1.000$ at every $k$, with
$\rho = 0.05$ for the unconstrained map and $0.83$--$0.92$ for the constrained
one. NL-DAS+r adds a reconstruction penalty ($\lambda_\text{recon} = 0.1$) and
reaches lower accuracy than NL-DAS on the same task. Single runs per cell.}
\label{tab:vacuity}
\end{table}

On indirect object identification the two nonlinear maps are indistinguishable by
accuracy and differ by more than an order of magnitude in $\rho$
(Table~\ref{tab:vacuity}). Adding a reconstruction penalty to the unconstrained
map lowers its accuracy by $0.13$ to $0.30$ on this task, which is consistent
with part of its accuracy depending on the absence of that penalty.

On modular arithmetic the pattern differs. Linear DAS also scores low at small
$k$ ($0.006$ at $k = 1$), so accuracy alone would favour the unconstrained
nonlinear map, which reaches $0.689$ at $k = 4$. Its $\rho$ over the same range is
$0.12$ to $0.43$. The constrained map reaches $1.000$ at $k = 1$ with
$\rho = 1.00$.

\subsection{Subspace dimension}
\label{sec:ksweep}

The unconstrained map's arithmetic accuracy could reflect an aligned space that
is too small, in which case accuracy should rise with $k$.

\begin{table}[t]
\centering
\small
\begin{tabular}{l rrrrrr}
\toprule
& $k{=}1$ & $k{=}2$ & $k{=}4$ & $k{=}8$ & $k{=}16$ & $k{=}32$ \\
\midrule
DAS (linear)          & 0.006 & 0.022 & 0.283 & 1.000 & 1.000 & 1.000 \\
NL-DAS                & 0.083 & 0.422 & 0.689 & 0.770 & 0.861 & 0.806 \\
\quad $\rho_\text{NL}$ & 0.12 & 0.29 & 0.43 & 0.46 & 0.54 & 0.33 \\
Structured VAE        & 1.000 & 1.000 & 1.000 & 1.000 & 1.000 & 1.000 \\
\quad $\rho_\text{VAE}$ & 1.00 & 1.00 & 1.00 & 1.00 & 0.99 & 1.00 \\
\bottomrule
\end{tabular}
\caption{Grokked modular addition (test accuracy $0.9996$), subspace dimension 1
to 32. Single runs per cell.}
\label{tab:ksweep}
\end{table}

Across a $32\times$ range the unconstrained map does not exceed $0.861$
(Table~\ref{tab:ksweep}), and both its accuracy and its diversity ratio decline
from $k = 16$ to $k = 32$. Linear DAS reaches $1.000$ at $k = 8$. Eight
dimensions is also the size of the span of four Fourier frequencies in two
components each, the representation modular arithmetic models of this kind are
reported to learn \citep{nanda2023progress}; we note the correspondence without
testing it. These are single runs, so the decline at $k = 32$ in particular
should not be read as an established effect.

\subsection{Language model tasks}
\label{sec:gpt2}

\begin{table}[t]
\centering
\small
\begin{tabular}{lcccc}
\toprule
& \multicolumn{2}{c}{DAS} & \multicolumn{2}{c}{Structured VAE} \\
\cmidrule(lr){2-3}\cmidrule(lr){4-5}
Task & Standard & Strict & Standard & Strict \\
\midrule
Indirect object identification & 0.19 & 0.19 & 1.00 & 1.00 \\
Gender bias                    & 0.62 & 0.52 & 1.00 & 1.00 \\
Greater-than                   & 0.93 & 0.28 & 1.00 & 1.00 \\
Subject--verb agreement        & 0.92 & 0.00 & 1.00 & 1.00 \\
Hypernymy                      & 0.07 & 0.02 & 0.97 & 0.90 \\
Capitals$^\dagger$             & 0.12 & 0.00 & 0.51 & 0.26 \\
\bottomrule
\end{tabular}
\caption{Six GPT-2 tasks at $k = 1$, layer 8, additive intervention. Standard
accuracy counts a success when the counterfactual answer outranks the base
answer; strict accuracy requires it to be the argmax over the vocabulary.
$^\dagger$Capitals has 182 classes over 190 examples.}
\label{tab:nlp}
\end{table}

The two scoring rules disagree substantially for linear DAS
(Table~\ref{tab:nlp}). On subject--verb agreement it scores $0.92$ standard and
$0.00$ strict, and on greater-than $0.93$ against $0.28$. Reporting only the
standard figure would overstate how often the counterfactual token is produced.
The structured VAE reaches strict accuracy at or above $0.90$ on five of six
tasks, and $0.26$ on capitals, where the number of classes approaches the number
of examples.

This table compares a linear map against a nonlinear one with more parameters, so
it does not isolate the effect of the constraints. We report it to indicate the
range of tasks on which the constrained map operates.

\subsection{Distribution shift}
\label{sec:shift}

\begin{table}[t]
\centering
\small
\setlength{\tabcolsep}{5pt}
\begin{tabular}{l rr rr rr}
\toprule
& \multicolumn{2}{c}{Standard} & \multicolumn{2}{c}{Held-out entities} & \multicolumn{2}{c}{Held-out templates} \\
\cmidrule(lr){2-3}\cmidrule(lr){4-5}\cmidrule(lr){6-7}
Method & IIA & $\rho$ & IIA & $\rho$ & IIA & $\rho$ \\
\midrule
DAS                       & 0.183 & ---  & 0.035 & ---  & 0.166 & ---  \\
NL-DAS                    & 1.000 & 0.05 & 1.000 & 0.05 & 0.537 & 0.05 \\
Structured VAE            & 1.000 & 0.83 & 0.001 & 0.78 & 0.963 & 0.51 \\
\quad + interchange loss  & 1.000 & 0.83 & 0.001 & 0.80 & 0.993 & 0.53 \\
\bottomrule
\end{tabular}
\caption{Indirect object identification at $k = 1$ under two distribution
shifts. Held-out entities withholds all test entities from training; held-out
templates withholds all test sentence structures.}
\label{tab:shift}
\end{table}

The two shifts separate the maps in opposite directions
(Table~\ref{tab:shift}). On held-out entities the unconstrained map retains
$\IIA = 1.000$ while the constrained map drops to $0.001$. On held-out templates
the constrained map reaches $0.993$ and the unconstrained map falls to $0.537$.
The unconstrained map's $\rho$ is $0.05$ on all three splits.

Reading which outcome is preferable requires an assumption about the task. The
indirect object identification variable is plausibly token-specific, in which
case transplanting it between disjoint entity sets has no well-defined target and
low accuracy is the appropriate result. Sentence structure is plausibly
irrelevant to it, in which case transfer across templates should succeed. Under
that reading the constrained map behaves as expected on both splits and the
unconstrained map on neither. We state the assumption rather than test it; a
reader who rejects it should take Table~\ref{tab:shift} as showing only that the
two maps respond differently to the two shifts.

\subsection{Which constraints matter}
\label{sec:ablation}

The structured VAE combines a label-conditional prior over the causal latent
\citep{khemakhem2020variational} with a supervised/unsupervised latent partition
\citep{narayanaswamy2017learning}, and expands the causal latent eightfold under
an $L_1$ penalty.

\begin{table}[t]
\centering
\small
\setlength{\tabcolsep}{5pt}
\begin{tabular}{lccc rrrr}
\toprule
Model & Prior & Partition & Expansion & Add. & Mult. & Quartic & IOI \\
\midrule
Plain VAE            & $\mathcal{N}(0,I)$ & \checkmark & --- & \TODO{--} & \TODO{--} & \TODO{--} & 0.68 \\
Label-conditional    & $p(z \mid y)$ & --- & --- & 0.02 & 0.00 & 0.02 & 0.56 \\
\quad + expansion    & $p(z \mid y)$ & --- & \checkmark & 1.00 & 0.72 & 0.32 & 0.83 \\
\quad + partition    & $p(z \mid y)$ & \checkmark & --- & 0.00 & 0.00 & 0.00 & 0.95 \\
\quad + both         & $p(z \mid y)$ & \checkmark & \checkmark & 1.00 & 1.00 & 1.00 & 0.98 \\
\bottomrule
\end{tabular}
\caption{Interchange accuracy at $k = 2$ by constraint combination. Single runs
per cell.}
\label{tab:ablation}
\end{table}

No single constraint reaches the accuracy of the full combination
(Table~\ref{tab:ablation}). The label-conditional prior alone reaches $0.02$ on
modular addition, and adding the partition to it reaches $0.00$. The expansion
alone reaches $1.00$ on addition, $0.72$ on multiplication, and $0.32$ on quartic
sum. The two together reach $1.00$ on all three. These are single runs, and we
have not tested whether the differences exceed run-to-run variation.

\TODO{SPARSITY CONFOUND, ABLATION RUNNING. Expansion and $L_1$ move together in
this table, because the non-expanded builders take no expansion factor. With
$\lambda = 10^{-3}$ the penalty is weak, so the expansion may account for the
whole effect. \texttt{experiments/sparsity\_vs\_overcompleteness\_ablation.py}
varies expansion $\in \{1,8\}$ against $\lambda \in \{0, 10^{-3}, 1\}$. Report
whichever result it returns, and remove the $L_1$ from the method description if
it proves inert.}

\subsection{Randomly initialised model}
\label{sec:randomnet}

\TODO{RUNNING. Both arms launched: GPT-2 with pretrained weights and with every
parameter re-initialised, indirect object identification, otherwise identical
settings. On the random arm the counterfactual target is defined by the task
rather than by the model's own outputs, since a random model is correct on
nothing. Report IIA and $\rho$ for NL-DAS and the structured VAE on both arms. If
the structured VAE also reaches high accuracy on the random arm, the reading
given in \S\ref{sec:vacuity} is not supported and the paper needs rewriting
around that.}

\section{Discussion}
\label{sec:discussion}

\citet{sutter2025nonlinear} established that unrestricted nonlinear causal
abstraction is uninformative and demonstrated it on randomly initialised models.
Our measurements are consistent with the restriction mattering: two maps
optimising the same interchange objective, differing in the presence of
reconstruction and regularisation terms, reach the same accuracy and differ in
diversity ratio and in behaviour under distribution shift. We have not shown that
the constrained map recovers the model's causal variable, and $\rho$ is a
statistic we introduce rather than an established measure. The strongest
supported claim is that accuracy alone did not separate these two maps while two
other measurements did.

The constraints involved are standard components of generative models
\citep{khemakhem2020variational, narayanaswamy2017learning}, and unsupervised
disentanglement is known to be unidentifiable without inductive biases of this
kind \citep{locatello2019challenging}. Nothing in the setup is specific to
interpretability, so the comparison could be repeated with other constrained
families.

For practitioners the implication is procedural. Interchange accuracy is
necessary for a causal variable claim and, on the evidence here, insufficient to
distinguish two maps that differ substantially by other measures. Reporting a
diversity ratio, performance on a randomly initialised model of the same
architecture, and behaviour on a split where transfer is not expected costs
little, and each of these separated the two maps in this study.

\paragraph{Limitations.}
Most cells in Tables~\ref{tab:ksweep} and~\ref{tab:ablation} are single runs, so
differences smaller than run-to-run variation cannot be distinguished, and we
have not measured that variation. The randomly initialised control in
\S\ref{sec:randomnet} is the experiment most directly comparable to
\citet{sutter2025nonlinear} and is not yet complete. The reading of the
held-out-entity result depends on an assumption about the task that we state and
do not test. The diversity ratio is a proxy: a map could preserve activation
spread without tracking the causal variable, and we do not rule this out.
Comparisons between linear DAS and the structured VAE are not capacity-controlled;
the constrained-against-unconstrained comparison is closer to matched, with the
constrained map using fewer parameters. Arithmetic experiments use synthetic
modular arithmetic with small transformers ($d_\text{model} = 128$). Reported
values come from separate runs rather than one sweep, which is why DAS on
indirect object identification appears at both $0.183$ and $0.19$.

These results describe two alignment maps on a small number of tasks. Whether the
separation holds for other constrained families, other tasks, or other models is
open.

\section{Methods}
\label{sec:methods}

\subsection{Interchange interventions}

Given a base input $b$ and a source input $s$, an interchange intervention
replaces part of the model's representation of $b$ with the corresponding part
from $s$ and asks whether the output matches the algorithm's counterfactual.
Writing $f$ for an encoder and $g$ for a decoder, the intervened activation is
$h' = g(\mathrm{swap}_k(f(h_b), f(h_s)))$, where $\mathrm{swap}_k$ replaces the
first $k$ coordinates. Interchange accuracy is the fraction of pairs whose
forward pass from $h'$ yields the counterfactual output. Standard accuracy counts
a success when the counterfactual answer outranks the base answer; strict
accuracy requires it to be the argmax over the vocabulary. DAS restricts $f$ to
an orthogonal rotation with $g = f^{-1}$.

\subsection{Alignment methods}

NL-DAS prepends a multilayer perceptron featuriser to a linear projection and
appends a multilayer perceptron inverse, trained end to end on interchange loss
with no reconstruction term. NL-DAS+r adds a reconstruction penalty
$\lambda_\text{recon}\lVert g(f(h)) - h \rVert^2$ with
$\lambda_\text{recon} = 0.1$.

\subsection{Structured variational autoencoder}

The encoder maps activations to two latent blocks: $\zc$, supervised through a
linear classifier and given a label-conditional Gaussian prior
$\mathcal{N}(\mu_y, \sigma_y^2 I)$ with per-label parameters learned as
embeddings \citep{khemakhem2020variational}; and $\zn$, unsupervised with a
standard normal prior \citep{narayanaswamy2017learning}. Training minimises
\begin{equation}
\mathcal{L} = \underbrace{\lVert g(f(h)) - h \rVert^2}_{\text{reconstruction}}
+ \underbrace{\KL_\text{causal} + \KL_\text{nuisance}}_{\text{regularisation}}
+ \alpha \underbrace{\mathcal{L}_\text{CE}(y, \hat{y})}_{\text{supervision}}
+ \lambda \underbrace{\lVert \zc \rVert_1}_{\text{sparsity}}.
\label{eq:loss}
\end{equation}
The causal block is expanded eightfold. All runs use $\zn$ dimension 16,
expansion factor 8, $\alpha = 10.0$, $\lambda = 10^{-3}$, learning rate
$10^{-3}$, and Adam. Encoder and decoder hidden width is 256 for GPT-2 tasks and
128 for modular arithmetic.

The label-conditional prior follows the identifiable VAE construction
\citep{khemakhem2020variational}. That construction's identifiability result
requires conditions we do not verify, including sufficient variability of the
per-label parameters relative to latent dimension, which is unlikely to hold on
our two-class tasks. We use the architecture as a source of constraints and make
no identifiability claim.

\subsection{Constrained interchange training}

After a reconstruction-and-classification warmup, an interchange term is added to
Equation~\ref{eq:loss} rather than substituted for it. Swapping $\zc$ and
decoding gives an intervened activation, which runs through the remaining layers
to contribute $-\log p(y_s \mid h')$. The generative terms stay active, so this
method optimises the same objective as NL-DAS with additional constraints.

The intervention is applied additively as
$h' = h_b + g(z_\text{swap}) - g(z_\text{orig})$, injecting the difference the
swap induces rather than the reconstructed activation, so that residual
reconstruction error cancels.

\subsection{Intervention diversity}

For each source label we compare the spread of intervened activations against the
spread of natural ones:
\begin{equation}
\rho = \frac{\mathbb{E}_{y_s}\bigl[\mathrm{std}(h'_{y_s})\bigr]}
            {\mathbb{E}_{y_s}\bigl[\mathrm{std}(h_{y_s})\bigr]},
\label{eq:rho}
\end{equation}
where the numerator ranges over intervened activations sharing source label $y_s$
and the denominator over the corresponding natural activations. Linear DAS
preserves the orthogonal complement by construction, so $\rho$ is reported only
for nonlinear maps.

\subsection{Tasks and controls}

Grokking experiments train one-layer transformers ($d_\text{model} = 128$) on
modular arithmetic at $p = 113$ to test accuracy above $0.999$. Language tasks use
GPT-2 at layer 8. Held-out entity and held-out template splits withhold,
respectively, all test entities and all test sentence structures from training.

Controls comprise random-label supervision, reconstruction-only training
($\alpha = 0$), an untrained probe, and a plain VAE without the latent partition.
Each is reported against the accuracy of a random subspace on the same task
rather than against zero, because chance accuracy varies with the number of
classes: on gender bias, with two classes, a random subspace and three of the
four controls all score $0.267$.

\TODO{CONTROL PROVENANCE. An earlier draft reported these controls as ``all below
$0.09$'', which is the figure from many-class tasks read as a universal bound.
Recompute per task against the random-subspace baseline already present in the
result files, and report $\rho$ for each.}

\section*{Data Availability}

\TODO{Repository and archived snapshot.}

\section*{Code Availability}

\TODO{Repository.}

\bibliographystyle{plainnat}
\bibliography{references,new_bib_entries}

\end{document}
